\scanpage{104}
\mspage{52}

$F$ comme ens. des sommets, et $A$ comme ens. d'arêtes [\deleted{\unsure{une arête est}}\inserted{graphe strict} : toute arête est non fermée, et déterminée par l'ens de ses extrémités] -- \deleted{ici} les sommets de ce graphe sont d'ordre $3$ (\deleted{\unsure{correspondant}}\inserted{les arêtes issues} du ``sommet''\ $f$ du graphe correspondent aux trois parties à deux éléments de $f\in\unsure{\mathscr P|\mathcal P|\mathfrak P}_3(S)$).

On pose

\[
\Pi=\Pi_0(\Gamma)\qquad\substack{\text{ens des composantes}\\\text{connexes du}\\\text{graphe }\Gamma}\tag*{(6)}
\]

\deleted{et pour $i\in\Pi$} de sorte qu'on a des applications surjectives

\[
\left\{\begin{aligned}
 F&\xrightarrow{\tau_F}\Pi\\
 A&\xrightarrow{\tau_A}\Pi
\end{aligned}\right.\tag*{(7)}
\]

dont les fibres en $i\in\Pi$ sont notées $F_i,A_i$, donc on a

\[
\left\{\begin{aligned}
 F&=\coprod_{i\in\Pi}F_i\qquad\text{les }F_i\ne\varnothing\\
 A&=\coprod_{i\in\Pi}A_i\qquad\text{les }A_i\ne\varnothing
\end{aligned}\right.\tag*{(8)}
\]

et

\[
\left\{\begin{aligned}
 &\left\{\begin{aligned}
 &\text{Si }\ a\in A,f\in F,\ a\subset f,\ \text{alors }\ f\in F_i\iff \deleted{\unsure{f}}a\in A_i\\
 &\text{Si }f,f'\in F_i,\ \text{il existe une chaîne }f_0=f,f_1,\ldots,f_N=f'
 \end{aligned}\right.\\
 &\text{de faces \deleted{mutuellement}\inserted{dont deux consécutives sont} adjacentes (i.e. }f_i\cap f_{i+1}\in A\text{ i.e}\\
 &\operatorname{card} f_i\cap f_{i+1}=2\text{) qui les joignent.}
\end{aligned}\right.\tag*{(9)}
\]
% archive sheet number: 103
